\section{Technical Preliminaries}

In this appendix, we review some technical preliminaries that will be useful for our analysis.  In \Cref{app:stoch_control}, we review some classical results from the theory of linear quadratic stochastic control, which we will use to derive our algorithm. Then in \Cref{app:ar_processes} we describe some classical results on autoregressive processes, which we will use to analyze the convergence of our algorithm in the quadratic optimization setting.



\subsection{Linear Quadratic Gaussian (LQG) Problems}\label{app:stoch_control}

In this section, we briefly review several classical results from the theory of linear quadratic stochastic control, which will be useful for our analysis.  Classical references for this material include \citet{kwakernaak1972linear,yong1999stochastic} among many others.

A \emph{Linear Quadratic Gaussian} (LQG) control problem is a stochastic control problem where the dynamics of the system are linear, the cost function is quadratic, and the noise is Gaussian.  More particularly, we suppose that $x_t^u \in \rr^d$ is a stochastic process satisfying
\begin{align}\label{eq:lqg-dynamics}
    d x_t^u = \left( \bA x_t^u + \bB u_t \right) d t + \bC W_t,
\end{align}
where $\bA, \bB, \bC$ are matrices of appropriate dimensions, $u_t$ is a control input (and thus progressively measurable with respect to the filtration generated by $W_t$), and $W_t$ is a standard Brownian motion in $\rr^d$.  For fixed time horizon $T$, we then seek to minimize a cost function of the form
\begin{align}\label{eq:lqg-cost}
    J(u) = \ee\left[ x_T^{u \top} \bQ_f x_T^u + \int_0^T \left( x_t^{u \top} \bQ x_t^u + u_t^\top \bR u_t \right) d t \right],
\end{align}
where $\bQ_f, \bQ, \bR$ are positive semidefinite matrices of appropriate dimensions.  The intuition for this cost is that we wish to minimize the expected deviation of the \emph{state} $x_t^u$ from the origin, while also ensuring that the control inputs $u_t$ are not too large; we pay both a \emph{final} cost for the end state $x_T^u$ being far from the origin as well as a \emph{running} cost for both intermediate states being far from the origin and control inputs being large.  This problem is classical and has been well-studied in the control theory literature.  We recall the fundamental result that the optimal control strategy is given by a linear state feedback controller, where the feedback gain is given by the solution to a matrix Riccati differential equation \citep{georgiou2013separation,kwakernaak1972linear}.
\begin{theorem}[\cite{yong1999stochastic} Theorem 6.1]\label{thm:lqg}
    Suppose that $x_t^u$ evolves according to \eqref{eq:lqg-dynamics} for some progressively measurable $u_t$.  Suppose further that $\bQ, \bQ_f, \bR$ are all symmetric and positive semidefinite and that $\bR$ is positive definite.  Let $\bP: [0, T] \to \rr^{d \times d}$ be the solution to the matrix Riccati differential equation
    \begin{align}\label{eq:riccati_general}
        - \bPdot_t = \bA^\top \bP_t + \bP_t \bA - \bP_t \bB \bR^{-1} \bB^\top \bP_t + \bQ, \quad \bP_T = \bQ_f,
    \end{align}
    where $\bPdot_t$ denotes the time derivative of $\bP_t$.  Then it holds that the optimal control $u_t$ that minimizes the cost function $J(u)$ defined in \eqref{eq:lqg-cost} is given by
    \begin{align}\label{eq:lqg-optimal-control}
        u_t = - \bR^{-1} \bB^\top \bP_t x_t^u
    \end{align}
    and the optimal cost is given by
    \begin{align}\label{eq:lqg-optimal-cost}
        J(u) = x_0^\top \bP_0 x_0 + \int_0^T \trace\left( \bP_t \bC \bC^\top \right) d t.
    \end{align}
\end{theorem}
While the Riccati equation provides the optimal controller in full generality, we will not require this in the sequel.  We thus state two corollaries, corresponding to special cases of the LQG problem, that will be sufficient for our purposes.  In the first, we will ignore any final state cost (so $\bQ_f = 0$), and assume that both $\bQ$ and $\bR$ are scalar multiples of the identity.
\begin{corollary}\label{cor:lqg-scalarcost}
    Suppose that $x_t^u$ evolves according to \eqref{eq:lqg-dynamics} for some progressively measurable $u_t$, and that $\bQ_f = 0$, $\bQ = q \cdot \eye$, and $\bR = r \cdot \eye$ for some $q, r > 0$.  Suppose further that $\bB = \eye$.  Then the optimal control $u_t$ that minimizes the cost function $J(u)$ defined in \eqref{eq:lqg-cost} is given by
    \begin{align}
        u_t = - r^{-1} \bP_t x_t^u, \quad \text{where} \quad - \bPdot_t = \bA^\top P_t + P_t \bA - r^{-1} P_t^2 + q \cdot \eye, \quad P_T = 0.
    \end{align}
\end{corollary}


In the second corollary, we will now ignore the \emph{running} state cost (so $\bQ = 0$), and assume that $\bQ_f$ and $\bR$ are scalar multiples of the identity.
\begin{corollary}\label{cor:lqg-finalcost}
    Suppose that $x_t^u$ evolves according to \eqref{eq:lqg-dynamics} for some progressively measurable $u_t$, and that $\bQ = 0$, $\bQ_f = q \cdot \eye$, and $\bR = r \cdot \eye$ for some $q, r > 0$.  Suppose further that $\bB = \eye$.  Then the optimal control $u_t$ that minimizes the cost function $J(u)$ defined in \eqref{eq:lqg-cost} is given by
    \begin{align}
        u_t = - r^{-1} \bP_t x_t^u, \quad \text{where} \quad  \bPdot_t = \bA^\top P_t + P_t \bA - r^{-1} P_t^2, \quad P_T = q \cdot \eye.
    \end{align}
\end{corollary}

Finally, we specialize to the fully scalar case, where $\bA$ is diagonal and prove here a closed-form solution to the Riccati equation.  While this calculation is standard, we include it here for completeness.
\begin{lemma}\label{lem:scalar-riccati-solution}
    Let $p: [0, T] \to \rr$ be the solution to the scalar Riccati equation
    \begin{align}
        - \dot{p}(t) = 2 a p(t) - r^{-1} p(t)^2 + q, \quad p(T) = q_f.
    \end{align}
    Let
    \begin{align}
        \beta_{\pm} = r ( a \pm \omega), \quad \text{where} \quad \omega = \sqrt{a^2 + r^{-1} q}.
    \end{align}
    Let
    \begin{align}
        \kappa(t) = \frac{q_f - \beta_+}{q_f - \beta_-} \cdot e^{- 2 \omega(T - t)}.
    \end{align}
    Then
    \begin{align}\label{eq:scalar-riccati-solution}
        p(t) = \frac{\beta_+ - \kappa(t) \beta_-}{1 - \kappa(t)}.
    \end{align}
\end{lemma}
\begin{proof}
    Note that by factoring we have
    \begin{align}
        \dot{p}(t) = r^{-1} \left( p(t) - \beta_+ \right)\left( p(t) - \beta_- \right).
    \end{align}
    Thus, when $\omega > 0$, we have by a standard separation of variables argument that
    \begin{align}
        \frac{p(t) - \beta_+}{p(t) - \beta_-} = \frac{q_f - \beta_+}{q_f - \beta_-} \cdot e^{- 2 \omega(T - t)}.
    \end{align}
    The result follows by rearranging.
\end{proof}






\subsection{Autoregressive Processes and the Yule-Walker Equations}\label{app:ar_processes}

We now review some classical results on autoregressive processes, which we will use to analyze the convergence of our algorithm in the quadratic optimization setting.  We refer to \citep{shumway2006time,brockwell2002introduction} for a more detailed review of this material.  We begin with a definition of autoregressive processes.
\begin{definition}\label{def:ar_process}
    We say a sequence of random variables $\{X_k\}_{k \geq 0}$ is an autoregressive process of order $p$ (AR($p$)) if it satisfies, for $k \geq p$,
    \begin{align}
        X_k = \sum_{i=1}^p \phi_i X_{k-i} + \epsilon_k,
    \end{align}
    for some $\phi_1, \dots, \phi_p \in \rr$ and some sequence of i.i.d. random variables $\{\epsilon_k\}_{k \geq 0}$ with mean zero and finite variance.  The \emph{characteristic polynomial} of the AR($p$) process is given by
    \begin{align}
        \phi(z) = 1 - \sum_{i=1}^p \phi_i z^i.
    \end{align}
    Moreover, we say that the process is \emph{causal} if it can be represented as
    \begin{align}
        X_k = \sum_{j = 0}^\infty \psi_j \epsilon_{k-j},
    \end{align}
    with the coefficients $\psi_j$ satisfying $\sum_{j=0}^\infty |\psi_j| < \infty$.
\end{definition}
Formally causality is necessary for us to ensure a limiting variance; fortunately there exists a simple characterization of causality in terms of the roots of the characteristic polynomial.
\begin{proposition}[Property 3.1 \cite{shumway2006time}]\label{prop:causality}
    An AR($p$) process is causal if and only if the roots of the characteristic polynomial $\phi(z)$ lie outside the unit circle in the complex plane.  In this case, $X_k$ has a well-defined limiting variance.
\end{proposition}
We now recall the Yule-Walker equations, which relate the coefficients $\phi_i$ of an AR($p$) process to the autocovariance function $\gamma(j) = \ee[X_k X_{k+j}]$.
\begin{theorem}[Section 3 \cite{shumway2006time}]\label{thm:yule-walker}
    Let $\{X_k\}_{k \geq 0}$ be a causal AR($p$) process with coefficients $\phi_1, \dots, \phi_p$ such that $\epsilon_k$ has variance $\sigma^2$. Then the autocovariance function $\gamma(j) = \ee[X_k X_{k+j}]$ satisfies the following equations for all $j \geq 0$:
    \begin{align}
        \gamma(j) = \sum_{i=1}^p \phi_i \gamma(j - i) + \sigma^2 \cdot \delta_{j, 0},
    \end{align}
    In particular, if $p = 2$, it holds that
    \begin{align}
        \gamma(1) &= \phi_1 \cdot \gamma(0) \\
        \gamma(2) &= \phi_1 \cdot \gamma(1) + \phi_2 \cdot \gamma(0) \\
        \gamma(0) &= \phi_1 \cdot \gamma(1) + \phi_2 \cdot \gamma(2) + \sigma^2.
    \end{align}
\end{theorem}
Finally, we recall the following result on the limiting variance of an AR($2$) process.
\begin{lemma}\label{lem:ar2-limiting-variance}
    Let $\{X_k\}_{k \geq 0}$ be a causal AR($2$) process with coefficients $\phi_1, \phi_2$ such that $\epsilon_k$ has variance $\sigma^2$ with fixed, deterministic $X_1, X_2$.  If $\phi(z) \neq 0$ for all complex $|z| \leq 1$, then the limiting variance of $X_k$ is given by
    \begin{align}
        \lim_{k \to \infty} \ee[X_k^2] = \gamma(0).
    \end{align}
\end{lemma}
\begin{proof}
    This follows immediately from Definition 3.7 and Property 3.1 of \citet{shumway2006time}. Indeed, let $X_k'$ denote the AR($2$) process with the same coefficients $\phi_1, \phi_2$ and noise $\epsilon_k$ but with zero initial conditions $X_1' = X_2' = 0$.  By causality, $X_k'$ has a well-defined limiting variance given by $\gamma(0)$.  Thus it suffices to show that $Y_k = X_k - X_k'$ has zero variance in the limit.  Observe that
    \begin{align}
        Y_k = \phi_1 Y_{k-1} + \phi_2 Y_{k-2},
    \end{align}
    so it remains to show that $Y_k \to 0$ as $k \to \infty$.  Writing out the transition matrix, we see that the eigenvalues thereof are the reciprocals of the roots of $\phi(z)$.  Because $\phi(z) \neq 0$ for all $|z| \leq 1$, the eigenvalues of the transition matrix are strictly less than $1$ in magnitude, so $Y_k \to 0$ as $k \to \infty$.
\end{proof}
